\section{Maldición de la Dimensionalidad y Complejidad}

% ================================================================
\begin{frame}{La Maldición de la Dimensionalidad (\textit{Curse of Dimensionality})}
\begin{columns}[T,onlytextwidth]
\begin{column}{0.50\textwidth}
\begin{contrastbox}{warn}{El Hipervolumen Vacío}
\scriptsize
En $[0, 1]^p$, capturar una fracción $r \in (0, 1)$ del volumen total exige una arista local:
\[
e_p(r) = r^{\frac{1}{p}}
\]
\begin{itemize}\setlength{\itemsep}{0.05em}
  \item $p = 1$: $e_1(0.10) = 0.10$ ($10\%$ del rango).
  \item $p = 10$: $e_{10}(0.10) \approx 0.80$ ($80\%$ del rango).
  \item $p = 100$: $e_{100}(0.10) \approx 0.98$ ($98\%$ del rango).
\end{itemize}
\end{contrastbox}
\end{column}

\begin{column}{0.48\textwidth}
\begin{ccdebox}
\scriptsize
\textbf{Consecuencia Teórica para KNN:}\\[0.05cm]
En alta dimensión, una vecindad de $k$ puntos \textbf{deja de ser local}: abarca casi todo el espacio muestral.
\end{ccdebox}

\vspace{0.06cm}
\begin{keybox}
\scriptsize
\textbf{Equidistancia Asintótica}:
\[
\lim_{p \to \infty} \frac{d_{\max} - d_{\min}}{d_{\min}} = 0
\]
Todos los puntos quedan aproximadamente a la misma distancia, diluyendo el concepto de cercanía.
\end{keybox}
\end{column}
\end{columns}
\source{Hastie, Tibshirani \& Friedman (2009), cap. 2; Donoho (2000).}
\end{frame}

% ================================================================
\begin{frame}{Estructuras de Indexación y Complejidad}
\begin{columns}[T,onlytextwidth]
\begin{column}{0.48\textwidth}
\begin{ccdebox}
\scriptsize
\textbf{Algoritmos de Búsqueda}
\begin{itemize}\setlength{\itemsep}{0.08em}
  \item \textbf{Fuerza Bruta (\texttt{brute})}: $O(N \cdot p)$ por consulta; evalúa todas las distancias.
  \item \textbf{KD-Tree (\texttt{kd\_tree})}: Partición ortogonal; eficiente con $p \le 20$ ($O(p \log N)$).
  \item \textbf{Ball-Tree (\texttt{ball\_tree})}: Hiperesferas anidadas; óptimo para métricas no euclidianas.
\end{itemize}
\end{ccdebox}
\end{column}

\begin{column}{0.48\textwidth}
\begin{contrastbox}{ccdemid}{Balance Metodológico}
\scriptsize
\textbf{Fortalezas}:
\begin{itemize}\setlength{\itemsep}{0.05em}
  \item Intuitivo y no paramétrico.
  \item Excelente \textit{baseline} no lineal.
  \item Extensión natural a multiclase.
\end{itemize}
\textbf{Limitaciones}:
\begin{itemize}\setlength{\itemsep}{0.05em}
  \item Latencia en inferencia ($O(N \cdot p)$).
  \item Sensible a variables irrelevantes.
  \item Almacenamiento continuo de datos.
\end{itemize}
\end{contrastbox}
\end{column}
\end{columns}
\source{Bentley (1975); Omohundro (1989); scikit-learn documentation.}
\end{frame}
